\chapter{Introduction}

Density-functional theory (DFT) rests on the Hohenberg--Kohn (HK) theorem: the ground-state density of an
electronic system determines the external potential up to an additive constant. In the continuum this was
proved rigorously only recently (via unique continuation). On finite lattices, which underlie the Hubbard
model, lattice DFT and every discretised calculation, the theorem is \emph{false} in general: Penz and van
Leeuwen exhibited lattice systems in which different potentials share a ground state
[J.\ Chem.\ Phys.\ 155, 244111 (2021)].

This blueprint describes a Lean~4/Mathlib formalization of lattice DFT. Its main result is that the
failures are exceptional:

\begin{quote}
\textbf{Generic Hohenberg--Kohn theorem} (Theorem~\ref{thm:generic_hk_lattice_fermions}).
For $N$ spinless fermions ($0<N<M$) on any connected graph with $M$ sites, hopping $t\neq 0$ and any
density--density interaction, the set of potentials $v\in\R^M$ for which some ground state of $H(v)$ is also
a ground state of a potential differing from $v$ by more than a constant is contained in a closed, nowhere
dense set of Lebesgue measure zero.
\end{quote}

The statement about potentials was left open by Penz and van Leeuwen [Quantum 7, 918 (2023)].

\section*{Structure of the proof}

\begin{enumerate}
  \item \textbf{Ground states and HK on configuration spaces} (Chapters~\ref{chap:ground_states},
  \ref{chap:hk}). Ground states are defined variationally. Weak HK (equal densities imply shared ground
  states) always holds. Strong HK can fail only if the ground state vanishes on some configurations.
  \item \textbf{Krylov criterion} (Chapter~\ref{chap:generic}). If an eigenvector vanishes on a
  configuration $c$, the Krylov matrix $[e_c, He_c, H^2e_c,\dots]$ is singular. Its determinant is a
  polynomial in $v$. So if all these \emph{Krylov polynomials} are nonzero (hypothesis (H$^*$)), HK fails
  only on the zero set of a nonzero polynomial, which is closed, nowhere dense and null.
  \item \textbf{Lattice fermions} (Chapter~\ref{chap:fermions}). Second quantisation is built from
  determinants, so fermionic signs are correct by construction; Slater determinants of single-particle
  eigenvectors are many-body eigenvectors.
  \item \textbf{Strong-potential limit} (Chapter~\ref{chap:strong}). Along $v=\lambda u$ the leading
  coefficient in $\lambda$ of a Krylov determinant does not depend on the interaction. It factorises as a
  Vandermonde determinant times diagonal coefficients $\alpha_y(y)$.
  \item \textbf{Path sums and resonance} (Chapter~\ref{chap:pathsum}). For free fermions,
  $\alpha_y(y)$ equals a signed sum over optimal particle assignments of products of single-particle route
  sums. Near a potential at which every moving particle is degenerate with its target, a single assignment
  dominates, so $\alpha_y(y)\not\equiv 0$.
  \item \textbf{Assembly} (Chapter~\ref{chap:main}).
\end{enumerate}

Each statement is linked to its Lean declaration; the dependency graph shows which results use which.
